\input{_preamble}
%% Hanging judgments inside an \altn stack (opt-in, with [phantomalign]).
%%
%% The motivating case, from a real document: two alternatives that differ
%% ONLY in their judgment.  Left-aligned, the stack put "est" under the
%% star and "sont" a star-width to its right, so the one pair of words the
%% reader is asked to compare was the one pair not aligned.  With alignment
%% on, the mark moves into a right-aligned gutter column of its own and the
%% real glyphs line up; the mark stays INSIDE the braces, because it judges
%% its own alternative and not the stack.
%%
%% Two things are pinned here and they fail independently.  The geometry is
%% measured on the ASCII probes, whose stems carry the sentinel names; the
%% French sentences are the original report, kept because they are what the
%% construction is actually for.  And the STARS THEMSELVES are counted:
%% before this, an alternative after the first lost its mark outright --
%% the tabular row separator \\ read it as its own starred form \\* -- so
%% "*sont" printed "sont", silently, with a clean compile.  That one is not
%% about alignment at all and is pinned again, unaligned, in altn.tex.
\GlossPhantomAlign
\begin{document}

% Probe: same stem either side of the judgment, so "aligned" is exactly
% "the two right edges agree" -- the left edge cannot tell the layouts
% apart, since a mark left inside its alternative starts at the column
% origin too.
\noindent PROBEL \lxAltn[l]{PSTEM}{*PSTEM} PROBER

% The yardstick for "no looser than typed".  The gutter butts the mark
% straight onto its word and adds no separation of its own, so the advance
% a hung mark occupies must equal the advance the same mark occupies in
% running text.  Both spellings of one stem are set here, so that advance
% can be had by subtraction without knowing anything about the font.
\noindent INLINE *QSTEM QSTEM ENDINLINE

% Two marks of different widths: the gutter sizes itself to the widest and
% every mark sits flush against the words, as \lx@hangjudge does at example
% level -- so the stems still share one left edge and "??" reaches further
% left than "*" does.
\noindent WIDEL \lxAltn[l]{*WSTEM}{??WSTEM}{WSTEM} WIDER

% No mark anywhere: the split layout must NOT engage, or a stack with
% nothing to align would gain the gutter gap and shift right under an
% option that is supposed to be about alignment only.  The same stack is
% set again with alignment OFF, and the two are compared to each other --
% measuring the gap from the sentinel's right edge to the first stem, which
% is brace plus both separations and is what a stray gutter would widen.
\noindent PLAINL \lxAltn[l]{NSTEM}{NSTEM} PLAINR

{\GlossPhantomAlignOff
\noindent OFFL \lxAltn[l]{OSTEM}{OSTEM} OFFR
\par}

% The original report.
\ex. \a. La secrétaire de Jean et collaboratrice de Paul
         \altn[l]{est}{*sont} à la gare.
\b. La secrétaire de Jean et la collaboratrice de Paul
    \altn[l]{*est}{sont} à la gare.
\z.
\end{document}
